\section{Sharp global deterministic theory under strong log-concavity}
\label{sec:global}

The global behavior separates into two mechanisms.  An underdispersed covariance must first reach
the curvature scale; after that, the objective localizes at a rate set by target conditioning.  For
the Fisher--Rao flow these mechanisms are visible in a single exact estimate, while for the two
discretizations they appear as a covariance bootstrap followed by geometric descent.  Throughout
this section we assume
\begin{equation}
\label{eq:global:curvature-assumption}
    \alpha\Id\preceq \nabla^2V(x)\preceq\beta\Id,
    \qquad x\in\R^d,
    \qquad
    \kappa:=\frac{\beta}{\alpha},
\end{equation}
and all Gaussian expectations are evaluated exactly.  We write
\begin{equation}
\label{eq:global:averaged-derivatives}
    \begin{aligned}
        g(a)&=-\E_{\N(m,C)}[\nabla V(X)],
        &
        A(a)&=\E_{\N(m,C)}[\nabla^2V(X)],\\
        R(a)&=C^{-1}-A(a).
    \end{aligned}
\end{equation}
The energy gap is
\(
\Delta(a)=\calE(a)-\calE(a_\star)
\), with \(\Delta_0=\Delta(a_0)\), and
\(\lambda_0:=\lambda_{\min}(C_0)\).

\subsection{Continuous covariance bootstrap and the intrinsic upper bound}
\label{sec:global:continuous}

The natural conditioning of the continuous flow is obtained by whitening with the covariance of
the optimizing Gaussian.  Set
\begin{equation}
\label{eq:global:optimizer-whitening}
    \widehat m=C_\star^{-1/2}(m-m_\star),
    \qquad
    \widehat C=C_\star^{-1/2}CC_\star^{-1/2},
    \qquad
    \widehat V(z)=V(m_\star+C_\star^{1/2}z).
\end{equation}
Then \(a_\star\) becomes \((0,\Id)\), and we define the optimal pointwise curvature constants
\begin{equation}
\label{eq:global:intrinsic-curvature}
    \alphastar
    :=\inf_z\lambda_{\min}(\nabla^2\widehat V(z)),
    \qquad
    \betastar
    :=\sup_z\lambda_{\max}(\nabla^2\widehat V(z)),
    \qquad
    \kappastar:=\frac{\betastar}{\alphastar},
\end{equation}
and the initial whitened covariance floor
\begin{equation}
\label{eq:global:intrinsic-initial-floor}
    \lambda_{0,\star}
    :=\lambda_{\min}(C_\star^{-1/2}C_0C_\star^{-1/2}).
\end{equation}
The first-order condition \(C_\star^{-1}=A(a_\star)\) gives
\(\alphastar\leq1\leq\betastar\).  Moreover,
\(1\leq\kappastar\leq\kappa^2\), although the upper comparison can be arbitrarily loose: an
anisotropic Gaussian target has \(\kappastar=1\).  Unlike \(\kappa\), all quantities in
\eqref{eq:global:intrinsic-curvature}--\eqref{eq:global:intrinsic-initial-floor} are unchanged by
invertible affine reparametrization.

The covariance floor follows most transparently from the precision.  If \(P_t=C_t^{-1}\), then
the Fisher--Rao covariance equation gives the exact Duhamel formula
\begin{equation}
\label{eq:global:precision-duhamel}
    \dot P_t=A(a_t)-P_t,
    \qquad
    P_t=e^{-t}P_0+\int_0^t e^{-(t-s)}A(a_s)\,\dd s.
\end{equation}
The same identity holds after optimizer whitening.

\begin{proposition}[Continuous covariance bootstrap]
\label{prop:global:continuous-bootstrap}
Along the Fisher--Rao flow,
\begin{equation}
\label{eq:global:continuous-floor}
    C_t\succeq
    \frac{\lambda_0e^t}{1+\beta\lambda_0(e^t-1)}\Id,
    \qquad
    \widehat C_t\succeq
    \ell_\star(t)\Id,
    \qquad
    \ell_\star(t):=
    \frac{\lambda_{0,\star}e^t}
         {1+\betastar\lambda_{0,\star}(e^t-1)}.
\end{equation}
Consequently, the two covariance floors are reached after the respective times
\begin{equation}
\label{eq:global:continuous-floor-times}
\begin{split}
    C_t\succeq(2\beta)^{-1}\Id
    &\quad\text{for }t\geq \logp\!\left(\frac{1}{\beta\lambda_0}\right),\\
    \widehat C_t\succeq(2\betastar)^{-1}\Id
    &\quad\text{for }t\geq
    \logp\!\left(\frac{1}{\betastar\lambda_{0,\star}}\right).
\end{split}
\end{equation}
\end{proposition}

\begin{proof}
Since \(A(a_t)\preceq\beta\Id\), \eqref{eq:global:precision-duhamel} implies
\(
P_t\preceq[e^{-t}\lambda_0^{-1}+\beta(1-e^{-t})]\Id
\); inversion gives the first bound.  Applying the same argument to the whitened flow gives the
second.  The two hitting times follow by making the decaying initial-precision term no larger than
the corresponding curvature term.
\end{proof}

\begin{theorem}[Affine-invariant continuous global rate]
\label{thm:global:continuous-upper}
For every \(t\geq0\),
\begin{equation}
\label{eq:global:continuous-decay}
    \boxed{
    \Delta(a_t)
    \leq
    \Delta_0
    \bigl[1+\betastar\lambda_{0,\star}(e^t-1)\bigr]^{-1/\kappastar}.}
\end{equation}
Thus, for \(0<\delta<\Delta_0\), the hitting time
\(
T_{\mathrm c}^{\glob}(\delta)
:=\inf\{t\geq0:\Delta(a_t)\leq\delta\}
\) satisfies
\begin{align}
\label{eq:global:continuous-hitting-exact}
    T_{\mathrm c}^{\glob}(\delta)
    &\leq
    \log\left(
        1+
        \frac{(\Delta_0/\delta)^{\kappastar}-1}
             {\betastar\lambda_{0,\star}}
    \right),\\
\label{eq:global:continuous-hitting-split}
    T_{\mathrm c}^{\glob}(\delta)
    &\leq
    \boxed{
    \logp\frac1{\betastar\lambda_{0,\star}}
    +\kappastar\log\frac{\Delta_0}{\delta}.}
\end{align}
The same argument in the original coordinates gives
\begin{equation}
\label{eq:global:continuous-hitting-original}
    T_{\mathrm c}^{\glob}(\delta)
    \leq
    \log\left(
        1+\frac{(\Delta_0/\delta)^\kappa-1}{\beta\lambda_0}
    \right);
\end{equation}
one may take the smaller of the intrinsic and original-coordinate bounds.
\end{theorem}

\begin{proof}
Optimizer whitening preserves both the energy and Fisher--Rao dissipation.  In whitened
coordinates the exact dissipation identity, the Gaussian Bures--Wasserstein PL inequality, and
\(\widehat C_t\succeq\ell_\star(t)\Id\) give
\begin{equation}
\label{eq:global:continuous-differential-inequality}
    \dot\Delta(a_t)
    =-\Gradnorm(a_t)^2
    \leq-\alphastar\ell_\star(t)\Delta(a_t).
\end{equation}
The scalar integral is explicit:
\begin{equation}
\label{eq:global:continuous-floor-integral}
    \int_0^t\ell_\star(s)\,\dd s
    =\frac1{\betastar}
      \log\bigl[1+\betastar\lambda_{0,\star}(e^t-1)\bigr].
\end{equation}
Gr\"onwall's inequality and \(\alphastar/\betastar=1/\kappastar\) prove
\eqref{eq:global:continuous-decay}.  Solving that estimate for \(t\) gives
\eqref{eq:global:continuous-hitting-exact}.  Finally, for \(R\geq1\) and \(b>0\),
\(
1+(R-1)/b\leq R/b
\) when \(b\leq1\), and it is at most \(R\) when \(b\geq1\); this yields
\eqref{eq:global:continuous-hitting-split}.
\end{proof}

The two terms in \eqref{eq:global:continuous-hitting-split} have different meanings.  The first is
the multiplicative covariance burn-in caused by underdispersion.  The second is the global
localization cost after the covariance reaches the intrinsic curvature scale.  The condition
number alone cannot encode arbitrary input size: even for a quadratic target, reaching a fixed
energy level from an unrestricted initial mean takes arbitrarily long.

\subsection{A matching continuous lower bound}
\label{sec:global:continuous-sharpness}

The intrinsic dependence in \Cref{thm:global:continuous-upper} cannot be replaced by a
dimension-free logarithmic dependence.  The obstruction already occurs in two dimensions and is
geometric: an anisotropic covariance can follow a slowly descending spiral valley while remaining
inside the globally admissible spectral band.

\begin{theorem}[Logarithmic-spiral lower bound]
\label{thm:global:spiral-lower}
For every \(K\geq16\), there is a globally \(C^3\) potential \(V_K:\R^2\to\R\) and an
initialization \(a_0=(m_0,C_0)\) such that
\begin{equation}
\label{eq:global:spiral-curvature}
    \alpha\Id\preceq\nabla^2V_K\preceq\beta\Id,
    \qquad
    \kappa=\frac\beta\alpha=2K+1,
    \qquad
    \frac1\beta\Id\preceq C_0\preceq\frac1\alpha\Id.
\end{equation}
Its optimizer-whitened condition number and initial covariance obey
\begin{equation}
\label{eq:global:spiral-intrinsic-data}
    K\leq\kappastar(V_K)\leq1.03K,
    \qquad
    \betastar\lambda_{0,\star}
    \geq\frac{1-\eta_K}{1+\eta_K},
    \qquad
    0\leq\eta_K\leq10^{-3}.
\end{equation}
For every level
\begin{equation}
\label{eq:global:spiral-local-level}
    \delta_{\loc}<\frac\alpha2e^{-0.08}\norm{m_0}^2,
\end{equation}
the exact Fisher--Rao flow satisfies
\begin{equation}
\label{eq:global:spiral-hitting-lower}
    \boxed{
    T_{\loc}:=\inf\{t\geq0:\Delta(a_t)\leq\delta_{\loc}\}
    \geq\frac K{100}
    \geq\frac{\kappastar(V_K)}{103}.}
\end{equation}
The construction extends to every dimension \(d\geq2\) by adjoining independent
quadratic coordinates initialized at equilibrium.
\end{theorem}

\begin{proof}[Proof outline]
The potential has a high-curvature logarithmic-spiral valley outside an isotropic quadratic core.
Its Hessian is certified globally between \(\alpha\Id\) and \(\beta\Id\).  A rescaled
four-dimensional system for radial displacement, angular displacement, and the two independent
covariance errors has a uniformly stable linear part; a semigroup estimate and a stopping-time
bootstrap control the nonlinear remainder and the Gaussian tails.  The resulting exact trajectory
obeys
\(
\norm{m_t}\geq e^{-0.04}\norm{m_0}
\) for \(0\leq t\leq K/100\).

The optimizing Gaussian is concentrated inside the quadratic core.  Hence
\(C_\star^{-1}\) is within relative error \(\eta_K\) of a scalar matrix, which gives
\eqref{eq:global:spiral-intrinsic-data}.  Strong convexity in the mean then yields
\(
\Delta(a_t)\geq(\alpha/2)\norm{m_t}^2
\), proving \eqref{eq:global:spiral-hitting-lower}.  The Hessian certificate, shadow system,
semigroup bounds, Gaussian-smoothing estimates, and dimension extension are given in
\Cref{app:spiral}.
\end{proof}

\begin{remark}[Scope of continuous sharpness]
\label{rem:global:spiral-scope}
The lower bound uses an admissible but anisotropic initial covariance.  It does not prove the same
bound under the stricter initialization \(C_0=\beta^{-1}\Id\).  The initial radius grows with
\(K\); if a prescribed positive Hessian-Lipschitz constant is imposed, it also grows with the
inverse of that constant.  These input-size factors belong to the logarithmic terms suppressed in
the shorthand statement
\(
T_{\loc}=\widetilde\Theta(\kappastar)
\).  No analogous linear lower bound is asserted for the endpoint of zero Hessian Lipschitz
constant, where the target is quadratic.
\end{remark}

\subsection{Covariance bootstraps and descent for the two discretizations}
\label{sec:global:discrete-upper}

Assume
\begin{equation}
\label{eq:global:initial-covariance-band}
    \lambda_{0,\min}\Id\preceq C_0\preceq\lambda_{0,\max}\Id,
    \qquad
    \lambda_0:=\lambda_{0,\min},
    \qquad
    \lambda_{\max}:=\max\left\{\lambda_{0,\max},\frac1\alpha\right\}.
\end{equation}
The upper barrier \(\lambda_{\max}\) controls the globally safe explicit stepsize.  The lower
barrier evolves much more favorably than a frozen estimate suggests.

\begin{proposition}[Two-sided covariance bootstrap]
\label{prop:global:discrete-bootstrap}
Let
\(
L_0=\min\{\lambda_0,(2\beta)^{-1}\}
\).

\emph{Riemannian retraction.}
If
\(
0<h\leq\min\{1,(\beta\lambda_{\max})^{-1}\}
\), then \(C_n\preceq\lambda_{\max}\Id\) and \(C_n\succeq L_n^{\Riem}\Id\), where
\begin{equation}
\label{eq:global:riemannian-floor-recurrence}
    L_{n+1}^{\Riem}
    =\frac{e^hL_n^{\Riem}}
           {1+(e-1)h\beta L_n^{\Riem}}.
\end{equation}
With
\begin{equation}
\label{eq:global:riemannian-floor-variables}
    y_0=\frac1{\beta L_0},
    \qquad
    y_\infty=\frac{(e-1)h}{e^h-1}\in[1,e-1),
\end{equation}
one has \(C_n\succeq(2\beta)^{-1}\Id\) for every
\begin{equation}
\label{eq:global:riemannian-floor-time}
    n\geq N_{\cov,-}^{\Riem}
    :=\left\lceil
        \frac1h\logp\frac{y_0-y_\infty}{2-y_\infty}
      \right\rceil
    =O\!\left(\frac1h\logp\frac1{\beta\lambda_0}\right).
\end{equation}
For the upper curvature band, put
\(
z_0=(\alpha\lambda_{0,\max})^{-1}
\) and
\(
z_\infty=h/(e^h-1)>1/2
\).  Then \(C_n\preceq2\alpha^{-1}\Id\) for all
\begin{equation}
\label{eq:global:riemannian-upper-band-time}
    n\geq N_{\cov,+}^{\Riem}:=
    \begin{cases}
        0, & z_0\geq\tfrac12,\\[3pt]
        \left\lceil
        \dfrac1h\log\dfrac{z_\infty-z_0}{z_\infty-1/2}
        \right\rceil, & z_0<\tfrac12.
    \end{cases}
\end{equation}

\emph{KL/Bregman resolvent.}
For every \(h>0\), \(C_n\preceq\lambda_{\max}\Id\) and
\(C_n\succeq L_n^{\KL}\Id\), where
\begin{equation}
\label{eq:global:kl-floor-recurrence}
    L_{n+1}^{\KL}
    =\frac{(1+h)L_n^{\KL}}{1+h\beta L_n^{\KL}},
    \qquad
    L_n^{\KL}
    =\frac1{\beta[1+(1/(\beta L_0)-1)(1+h)^{-n}]}.
\end{equation}
Consequently, \(C_n\succeq(2\beta)^{-1}\Id\) for
\begin{equation}
\label{eq:global:kl-floor-time}
    n\geq N_{\cov,-}^{\KL}
    :=\left\lceil
      \frac{\logp(1/(\beta L_0)-1)}{\log(1+h)}
    \right\rceil.
\end{equation}
When \(h\leq1\), this is
\(
O(h^{-1}\logp(1/(\beta\lambda_0)))
\).  In addition, \(C_n\preceq2\alpha^{-1}\Id\) for
\begin{equation}
\label{eq:global:kl-upper-band-time}
    n\geq N_{\cov,+}^{\KL}
    :=\left\lceil\frac{\log2}{\log(1+h)}\right\rceil.
\end{equation}
\end{proposition}

\begin{proof}
For the Riemannian update, let \(B_n=C_n^{1/2}A(a_n)C_n^{1/2}\).  The precision update is
\(
C_{n+1}^{-1}=e^{-h}C_n^{-1/2}e^{hB_n}C_n^{-1/2}
\).  The upper barrier follows from \(e^X\succeq\Id+X\).  Under
\(0\preceq hB_n\preceq\Id\), the chord bound
\(e^{hB_n}\preceq\Id+(e-1)hB_n\) gives
\eqref{eq:global:riemannian-floor-recurrence}.  Inverting the scalar recurrence produces
\(
(\beta L_n^{\Riem})^{-1}
=y_\infty+(y_0-y_\infty)e^{-nh}
\), which yields \eqref{eq:global:riemannian-floor-time}.  A lower affine recursion for
\(\lambda_{\min}(C_n^{-1})/\alpha\) gives
\eqref{eq:global:riemannian-upper-band-time}.

For the KL update the precision recursion is exactly affine:
\(
C_{n+1}^{-1}=(C_n^{-1}+hA(a_n))/(1+h)
\).  Comparing it with \(\alpha\Id\) and \(\beta\Id\) gives
\eqref{eq:global:kl-floor-recurrence} and the upper barrier.  Solving the resulting scalar
recurrences proves \eqref{eq:global:kl-floor-time}--\eqref{eq:global:kl-upper-band-time}.
\end{proof}

The descent mechanisms differ even though both methods share the mean update.  Let
\(
K_{\glob}:=\beta\lambda_{\max}\geq1
\), and suppose at the current iterate that
\(L\Id\preceq C\preceq\lambda_{\max}\Id\).  We use the Gaussian
Bures--Wasserstein gradient norm
\begin{equation}
\label{eq:global:wasserstein-gradient-norm}
    \norm{\grad^{W_2}\calE(a)}_{W_2}^2
    :=\norm{g(a)}^2+\Tr(R(a)C R(a)).
\end{equation}

\begin{proposition}[One-step deterministic descent]
\label{prop:global:one-step-descent}
For the Riemannian retraction, if \(0<h\leq K_{\glob}^{-1}\), then
\begin{equation}
\label{eq:global:riemannian-descent}
    \calE(a^+)
    \leq\calE(a)-h(1-K_{\glob}h)\Gradnorm(a)^2.
\end{equation}
In particular, \(h\leq(2K_{\glob})^{-1}\) gives decrease by at least
\((h/2)\Gradnorm(a)^2\).

For the KL/Bregman resolvent, if \(0<h\leq K_{\glob}^{-1}\), then the useful orientation of the
Gaussian divergence is forward:
\begin{align}
\label{eq:global:kl-forward-descent}
    \calE(a^+)
    &\leq\calE(a)-\frac1h\KL(\rho_{a^+}\,\|\,\rho_a),\\
\label{eq:global:kl-forward-control}
    \KL(\rho_{a^+}\,\|\,\rho_a)
    &\geq
    \frac{Lh^2}{4(1+hK_{\glob})^2}
    \norm{\grad^{W_2}\calE(a)}_{W_2}^2.
\end{align}
The reverse orientation \(\KL(\rho_a\|\rho_{a^+})\) does not provide the corresponding bound for
arbitrary mean displacement.
\end{proposition}

\begin{proof}[Proof idea]
For the Riemannian map, couple the old and new Gaussians through a shared standard normal variable.
Smoothness of \(V\), Gaussian integration by parts, and the scalar inequality
\(
e^w\leq1+w+w^2
\) for \(|w|\leq1\) control the exponential covariance retraction and yield
\eqref{eq:global:riemannian-descent}.

For the KL map, diagonalize
\(
C^{1/2}A(a)C^{1/2}
\).  Its covariance update is a rational function of this matrix, so the covariance part reduces
to a scalar inequality on eigenvalues; the mean part is controlled by \(hK_{\glob}\leq1\).
This gives \eqref{eq:global:kl-forward-descent}.  The elementary inequality
\(
q^{-1}-1+\log q\geq(q-1)^2/[2\max\{q,1\}^2]
\) then proves \eqref{eq:global:kl-forward-control}.
\end{proof}

\begin{theorem}[Global complexity of the deterministic discretizations]
\label{thm:global:discrete-upper}
Let
\(
N_\bullet^{\glob}(\delta)
:=\inf\{n\geq0:\Delta(a_n)\leq\delta\}
\) for \(0<\delta<\Delta_0\).

If the Riemannian retraction is run with
\begin{equation}
\label{eq:global:riemannian-global-step}
    0<h\leq\frac1{2\beta\lambda_{\max}},
\end{equation}
then
\begin{equation}
\label{eq:global:riemannian-global-complexity}
    \boxed{
    N_{\Riem}^{\glob}(\delta)
    =O\!\left(
        \frac1h\logp\frac1{\beta\lambda_0}
        +\frac\kappa h\logp\frac{\Delta_0}{\delta}
    \right).}
\end{equation}
If the KL/Bregman resolvent is run with
\begin{equation}
\label{eq:global:kl-global-step}
    0<h\leq\frac1{\beta\lambda_{\max}},
\end{equation}
then
\begin{equation}
\label{eq:global:kl-global-complexity}
    \boxed{
    N_{\KL}^{\glob}(\delta)
    =O\!\left(
        \frac1h\logp\frac1{\beta\lambda_0}
        +\frac\kappa h\logp\frac{\Delta_0}{\delta}
    \right).}
\end{equation}
The constants hidden by \(O(\cdot)\) are independent of the dimension and the conditioning
parameters.
\end{theorem}

\begin{proof}
After the covariance burn-in established in
\Cref{prop:global:discrete-bootstrap}, one has
\(C_n\succeq(2\beta)^{-1}\Id\).  The Fisher--Rao/Bures comparison and the Gaussian
Bures--Wasserstein PL inequality imply
\begin{equation}
\label{eq:global:discrete-pl-after-burnin}
    \Gradnorm(a_n)^2\geq\frac1{2\kappa}\Delta(a_n).
\end{equation}
For the Riemannian scheme,
\eqref{eq:global:riemannian-descent} therefore gives
\(
\Delta_{n+1}\leq(1-h/(4\kappa))\Delta_n
\).  For the KL scheme,
\eqref{eq:global:kl-forward-descent}--\eqref{eq:global:kl-forward-control} and
\(1+hK_{\glob}\leq2\) give
\(
\Delta_{n+1}\leq(1-h/(16\kappa))\Delta_n
\).  Adding the covariance burn-in and summing these geometric contractions proves
\eqref{eq:global:riemannian-global-complexity} and
\eqref{eq:global:kl-global-complexity}.
\end{proof}

The KL covariance recursion is an affine precision resolvent and preserves the global upper
barrier for every \(h>0\); the global descent theorem certifies \(h\leq
(\beta\lambda_{\max})^{-1}\).  The Riemannian method uses a matrix exponential and its stated
descent rate uses the smaller range \(h\leq(2\beta\lambda_{\max})^{-1}\).  These are certified
theoretical ranges, rather than a claim that either method is stable beyond them on every target.

\subsection{Fixed-step sharpness in one dimension}
\label{sec:global:discrete-sharpness}

At the globally safe scale \(h=\Theta(\kappa^{-1})\), the localization term in
\Cref{thm:global:discrete-upper} is \(\Theta(\kappa^2)\) for a constant-factor energy reduction.
This polynomial dependence is attained even when the covariance starts at the curvature scale and
the Hessian is Lipschitz with a constant independent of \(\kappa\).

\begin{theorem}[Bump-train lower bound for both fixed-step maps]
\label{thm:global:bump-lower}
Fix \(\gamma\in(0,1/2]\) and \(L_0>0\).  There exist constants
\(T,q,Y>0\) and \(\kappa_0\geq1\), depending only on \(\gamma\), \(L_0\), and a fixed smooth
bump profile, such that for every \(\kappa\geq\kappa_0\) there is an even potential
\(V_\kappa\in C^\infty(\R)\) satisfying
\begin{equation}
\label{eq:global:bump-regularity}
    1\leq V_\kappa''\leq\kappa,
    \qquad
    \norm{V_\kappa'''}_\infty\leq L_0,
    \qquad
    \kappastar(V_\kappa)=\kappa.
\end{equation}
Initialize
\begin{equation}
\label{eq:global:bump-initialization}
    m_0=\frac{Y\kappa^4}{\gamma},
    \qquad
    C_0=c_0=\frac1\kappa,
\end{equation}
and run either deterministic discretization with the fixed step
\begin{equation}
\label{eq:global:bump-step}
    h=\frac\gamma\kappa.
\end{equation}
Then
\begin{equation}
\label{eq:global:bump-gap-lower}
    \Delta_n\geq q\Delta_0
    \qquad
    \text{for every }0\leq n\leq
    \left\lfloor\frac{T\kappa^2}{\gamma}\right\rfloor.
\end{equation}
Consequently, for a constant-factor objective reduction,
\begin{equation}
\label{eq:global:bump-complexity-lower}
    N_{\Riem}=\Omega\!\left(\frac{\kappa^2}{\gamma}\right)
    =\Omega\!\left(\frac\kappa h\right),
    \qquad
    N_{\KL}=\Omega\!\left(\frac{\kappa^2}{\gamma}\right)
    =\Omega\!\left(\frac\kappa h\right).
\end{equation}
\end{theorem}

\begin{proof}[Proof outline]
The construction places flat-top curvature bumps along a prescribed scalar reference orbit.  Each
bump has height \(\kappa-1\) and width \(\Theta(\kappa)\), which keeps
\(\norm{V_\kappa'''}_\infty\) bounded.  At the matched covariance \(c_\kappa=1/\kappa\), the
Gaussian-averaged curvature satisfies
\(
c_\kappa A(m,c_\kappa)=1+o(1)
\), so both covariance maps are almost stationary.  Meanwhile
\(
b(m,c_\kappa)/m=\Theta(1)
\), and hence the relative mean displacement per step is only
\(
h c_\kappa b(m,c_\kappa)/m=\Theta(\gamma/\kappa^2)
\).

Gaussian leakage between neighboring bumps is small enough to shadow the reference orbit for
\(\Theta(\kappa^2/\gamma)\) steps.  Strong convexity in the mean and control of the shadowing error
then show that the objective gap retains a fixed fraction of its initial value.  The bump
placement, uniform regularity certificate, simultaneous shadowing of the two covariance maps, and
objective comparison are proved in \Cref{app:bump}.
\end{proof}

\begin{remark}[Scope of discrete sharpness]
\label{rem:global:bump-scope}
The target in \Cref{thm:global:bump-lower} may depend on the fixed value of \(\gamma\).  The theorem
proves sharpness of the factor \(\kappa/h\) for a constant-factor reduction by an explicit method
committed to one globally fixed step.  It does not cover adaptive relative-curvature steps,
backtracking or Armijo line searches, or implicit and Rosenbrock-stabilized updates, and it does not
claim that the entire factor \(\log(\Delta_0/\delta)\) is necessary at every target accuracy.
\end{remark}

\subsection{Removing the covariance burn-in}
\label{sec:global:warm-starts}

The burn-in can be removed without changing the subsequent Fisher--Rao dynamics.  Two optional
initializers are useful in different oracle models: a short Bures--Wasserstein transport stage uses
the averaged quantities \(g(a)\) and \(A(a)\), whereas a quadratic rescue uses one pointwise
gradient--Hessian query.

\begin{proposition}[Bures--Wasserstein transport warm start]
\label{prop:global:bures-warm-start}
Consider first the Gaussian Bures--Wasserstein flow
\begin{equation}
\label{eq:global:bures-flow}
    \dot m_t=g(a_t),
    \qquad
    \dot C_t=2\Id-C_tA(a_t)-A(a_t)C_t.
\end{equation}
If \(C_0\succeq\lambda_0\Id\), then
\begin{equation}
\label{eq:global:bures-floor}
    C_t\succeq
    \left[
        \frac1\beta+\left(\lambda_0-\frac1\beta\right)e^{-2\beta t}
    \right]\Id.
\end{equation}
For \(c\in(0,1)\), define
\begin{equation}
\label{eq:global:bures-warm-time}
    T_W(c):=
    \begin{cases}
        0, & \lambda_0\geq c/\beta,\\[3pt]
        \dfrac1{2\beta}
        \log\dfrac{1-\beta\lambda_0}{1-c},
        & \lambda_0<c/\beta.
    \end{cases}
\end{equation}
Then \(C_{T_W(c)}\succeq(c/\beta)\Id\), the energy has not increased, and a permanent switch to
the Fisher--Rao flow reaches \(\Delta\leq\delta\) within the additional time
\begin{equation}
\label{eq:global:bures-to-fr-continuous}
    \min\left\{
        \frac\kappa c\logp\frac{\Delta_0}{\delta},
        \logp\frac\kappa{c\betastar}
        +\kappastar\logp\frac{\Delta_0}{\delta}
    \right\}.
\end{equation}
For fixed \(c\), the warm-up \eqref{eq:global:bures-warm-time} is \(O(1/\beta)\), independently of
how small \(\lambda_0\) is.

There is also a one-step discrete version.  For any \(c\in(0,1]\) and \(a=(m,C)\), let
\begin{equation}
\label{eq:global:bures-forward-backward}
    \eta=\frac c\beta,
    \quad
    M=\Id-\eta A(a),
    \quad
    m^+=m+\eta g(a),
    \quad
    \widetilde C=MCM,
\end{equation}
and set
\begin{equation}
\label{eq:global:bures-forward-backward-covariance}
    C^+
    =\frac12\left(
        \widetilde C+2\eta\Id
        +[\widetilde C(\widetilde C+4\eta\Id)]^{1/2}
      \right).
\end{equation}
If \(C\preceq\lambda_{\max}\Id\), then
\(\eta\Id\preceq C^+\preceq\lambda_{\max}\Id\) and
\(\calE(a^+)\leq\calE(a)\).  Thus, if \(\lambda_0<c/\beta\), one such step followed by either
Fisher--Rao discretization replaces the burn-in term by one iteration and gives
\begin{equation}
\label{eq:global:bures-to-fr-discrete}
    N_{W\to\bullet}^{\glob}(\delta;c)
    =1+O\!\left(
        \frac\kappa{ch}\logp\frac{\Delta_0}{\delta}
    \right),
    \qquad \bullet\in\{\Riem,\KL\},
\end{equation}
under the corresponding stepsize restriction in
\eqref{eq:global:riemannian-global-step} or \eqref{eq:global:kl-global-step}.  If the initial floor
already exceeds \(c/\beta\), the transport step is skipped and the additive one is absent.
\end{proposition}

\begin{proof}[Proof idea]
An eigenvalue \(\mu_i(t)\) of \(C_t\) along \eqref{eq:global:bures-flow} satisfies
\(
\dot\mu_i=2(1-\mu_i h_i)
\) for some \(h_i\in[\alpha,\beta]\).  Scalar comparison with
\(
\dot\ell=2(1-\beta\ell)
\) proves \eqref{eq:global:bures-floor}.  The flow is the Gaussian Bures--Wasserstein gradient flow
of \(\calE\), so its energy decreases \citep{lambert2022variational}.  After the switch, the floor
\(c/\beta\) persists under Fisher--Rao evolution; the original-coordinate PL estimate gives the
first term in \eqref{eq:global:bures-to-fr-continuous}, while
\Cref{thm:global:continuous-upper} applied to
\(
\widehat C\succeq(c/\kappa)\Id
\) gives the second.

For the discrete step, the scalar map
\(
\psi_\eta(x)=(x+2\eta+\sqrt{x(x+4\eta)})/2
\) is increasing and satisfies \(\psi_\eta(0)=\eta\), proving the lower floor.  Its identity
\(
\psi_\eta((1-\eta/L)^2L)=L
\) proves the upper barrier.  Energy decrease is the standard Bures--Wasserstein
forward--backward inequality for \(\eta\leq1/\beta\) \citep{diao2023forward}.  The bounds in
\Cref{thm:global:discrete-upper}, restarted from the new floor, yield
\eqref{eq:global:bures-to-fr-discrete}.
\end{proof}

The additive covariance growth in \eqref{eq:global:bures-floor} is coordinate-dependent:
the Bures--Wasserstein flow is not affine invariant.  We use it only as a finite warm start
followed by a permanent switch; no general Wasserstein--Fisher--Rao dynamics is needed.

\begin{proposition}[One-query quadratic rescue]
\label{prop:global:quadratic-rescue}
Suppose a pointwise Hessian oracle is available.  Given \(m_{\mathrm{in}}\), query
\begin{equation}
\label{eq:global:quadratic-rescue-query}
    G_{\mathrm{in}}=\nabla V(m_{\mathrm{in}}),
    \qquad
    H_{\mathrm{in}}=\nabla^2V(m_{\mathrm{in}}),
\end{equation}
and initialize
\begin{equation}
\label{eq:global:quadratic-rescue}
    \boxed{
    m_0=m_{\mathrm{in}}-H_{\mathrm{in}}^{-1}G_{\mathrm{in}},
    \qquad
    C_0=H_{\mathrm{in}}^{-1}.}
\end{equation}
Then
\begin{equation}
\label{eq:global:quadratic-rescue-band}
    \frac1\beta\Id\preceq C_0\preceq\frac1\alpha\Id,
\end{equation}
so both deterministic Fisher--Rao schemes start with zero covariance burn-in.  In particular, with
\(\Delta_0\) denoting the post-rescue gap,
\begin{equation}
\label{eq:global:quadratic-rescue-complexity}
    N_{\mathrm{rescue}\to\bullet}^{\glob}(\delta)
    =1+O\!\left(
        \frac\kappa h\logp\frac{\Delta_0}{\delta}
    \right),
    \qquad
    \begin{cases}
        0<h\leq(2\kappa)^{-1}, & \bullet=\Riem,\\
        0<h\leq\kappa^{-1}, & \bullet=\KL,
    \end{cases}
\end{equation}
where the additive one counts the gradient--Hessian query.

If
\(
V(x)=\tfrac12(x-\mu)^\top H(x-\mu)+c
\) with \(H\succ0\), then \eqref{eq:global:quadratic-rescue} gives
\(
(m_0,C_0)=(\mu,H^{-1})=a_\star
\) and \(\Delta_0=0\).  Over the class of unknown Gaussian targets, the pointwise
gradient--Hessian oracle complexity of exact recovery is therefore exactly one.
\end{proposition}

\begin{proof}
The inverse map is operator antitone, so
\eqref{eq:global:curvature-assumption} immediately gives
\eqref{eq:global:quadratic-rescue-band}.  Hence
\(\lambda_0\geq1/\beta\), \(\lambda_{\max}=1/\alpha\), and
\Cref{thm:global:discrete-upper} starts directly in its localization phase, proving
\eqref{eq:global:quadratic-rescue-complexity}.  For a quadratic potential,
\(
G_{\mathrm{in}}=H(m_{\mathrm{in}}-\mu)
\) and \(H_{\mathrm{in}}=H\), which proves exact recovery.  Zero queries cannot recover every
unknown pair \((\mu,H)\), while one query returns both \(H\) and \(H(m_{\mathrm{in}}-\mu)\); thus one
query is also necessary.
\end{proof}

The transport warm start and quadratic rescue modify only the initialization stage.  They do not
remove the global-conditioning obstruction: the bump-train lower bound begins at
\(C_0=\kappa^{-1}\), after covariance burn-in has already ended.  Nor do they change the local
non-Gaussian stiffness studied in the next section.  Covariance warm-up, global localization, and
local contraction are therefore distinct parts of the deterministic rate landscape.
